\documentclass[10pt,aspectratio=169]{beamer}
\usepackage[utf8]{inputenc}
\usepackage[T1]{fontenc}
\usepackage[french]{babel}
\usepackage{amsmath}
\usepackage{amssymb}
\usepackage{booktabs}
\usepackage{array}
\usepackage{tikz}
\usetikzlibrary{positioning,calc,arrows.meta}
\usepackage{eurosym}

% Thème : celui des autres séances
\usetheme{Madrid}
\usecolortheme{default}

\definecolor{bleuactu}{RGB}{0,51,102}
\definecolor{vertactu}{RGB}{0,102,51}
\definecolor{rougeactu}{RGB}{153,0,0}
\definecolor{orangeactu}{RGB}{230,126,34}

\setbeamercolor{structure}{fg=bleuactu}
\setbeamercolor{frametitle}{bg=bleuactu,fg=white}

\newcommand{\highlight}[1]{\textbf{\color{rougeactu}#1}}
\newcommand{\BE}{\mathrm{BE}}
\newcommand{\MR}{\mathrm{MR}}
\newcommand{\SCR}{\mathrm{SCR}}
\newcommand{\keur}{\,k\euro}

% Compilation pdfLaTeX sans erreur : caractères absents de l'encodage T1 (voir README « LaTeX »)
\DeclareUnicodeCharacter{2248}{\ensuremath{\approx}}
\DeclareUnicodeCharacter{2264}{\ensuremath{\leq}}
\DeclareUnicodeCharacter{2265}{\ensuremath{\geq}}

\title[Séance 1 -- Pourquoi la comptabilité prudentielle ?]{Séance 1 : Pourquoi la comptabilité prudentielle ?}
\subtitle{Comptabilité prudentielle en assurance}
\author[F. Borel-Mathurin]{Fabrice BOREL-MATHURIN}
\institute[ACPR-DEAR et SAF Lyon 1]{
    ACPR - DEAR\\
    et\\
    Laboratoire SAF Université Claude Bernard Lyon 1
}
\date{2026-2027}

\begin{document}

\frame{\titlepage}

\begin{frame}{Plan de la séance (4 h)}
\tableofcontents
\end{frame}

\begin{frame}{Objectifs pédagogiques}
\begin{block}{À l'issue de cette séance, vous serez capable de :}
\begin{enumerate}
  \item Expliquer la finalité de la comptabilité prudentielle
  \item Identifier ce qui la distingue de la comptabilité financière
  \item Situer les institutions (EIOPA, ACPR) et les textes fondateurs
  \item Construire un bilan prudentiel simplifié à partir d'un bilan comptable
  \item Calculer et interpréter un ratio de couverture du SCR
\end{enumerate}
\end{block}
\vspace{0.3cm}
\begin{alertblock}{Déroulé}
\textbf{Cours :} 2 h \quad | \quad \textbf{Travaux dirigés :} 2 h sur tableur (TD 1 : le mini-bilan prudentiel de l'Assureur Bêta)
\end{alertblock}
\end{frame}

\begin{frame}{Pour commencer : trois questions}
\begin{block}{À discuter à deux, 5 minutes}
\begin{enumerate}
  \item Une obligation achetée 100\,\euro{} cote aujourd'hui 85\,\euro. Que vaut-elle dans les
        comptes de l'assureur ? Et pour quelqu'un qui voudrait la racheter ?
  \item Un assureur encaisse la prime \emph{avant} de payer les sinistres. Qu'est-ce que cela
        change pour celui qui le contrôle ?
  \item Qui a intérêt à ce qu'un assureur reste solvable ?
\end{enumerate}
\end{block}
\vspace{0.2cm}
\small Nous reviendrons sur vos réponses au fil de la séance.
\end{frame}

% =====================================================================
\section{Pourquoi une comptabilité prudentielle ?}
% =====================================================================

\begin{frame}{Le cycle de production inversé de l'assurance}
\begin{columns}[T]
\begin{column}{0.48\textwidth}
\begin{block}{Ce qui fait la spécificité de l'assureur}
\begin{itemize}
  \item Il \textbf{encaisse d'abord} la prime, il \textbf{paie ensuite}, parfois des décennies plus tard
  \item Le coût de revient n'est connu qu'\textbf{a posteriori}
  \item Ses engagements sont \textbf{longs} et \textbf{aléatoires}
\end{itemize}
\end{block}
\end{column}
\begin{column}{0.48\textwidth}
\begin{alertblock}{Conséquence}
Le résultat de l'année ne dit pas si l'assureur pourra payer ses assurés \textbf{même en cas de
choc}. Il faut une mesure \textbf{prospective} de sa solidité.
\end{alertblock}
\end{column}
\end{columns}
\vspace{0.4cm}
\begin{center}
\fbox{\parbox{0.8\textwidth}{\centering
Fonds propres = valeur des actifs $-$ valeur des engagements\\[1mm]
\small\emph{les deux évalués à leur valeur économique (directive 2009/138/CE, art.~75)}}}
\end{center}
\end{frame}

\begin{frame}{Comptabilité financière et comptabilité prudentielle}
\begin{table}
\centering
\small
\renewcommand{\arraystretch}{1.4}
\begin{tabular}{p{3.2cm}p{4.8cm}p{4.8cm}}
\toprule
 & \textbf{Comptabilité financière} & \textbf{Comptabilité prudentielle} \\
\midrule
Question posée & Quelle performance ? & L'assureur peut-il tenir ses engagements ? \\
Destinataires & Actionnaires, marché & Superviseur, assurés \\
Valorisation & Coût historique, principe de prudence & Valeur économique (« cohérente avec le marché ») \\
Horizon & L'exercice écoulé & Les engagements jusqu'à leur extinction \\
Indicateur clé & Résultat net & Ratio de couverture du SCR \\
Référentiels & Code de commerce, règlement ANC 2015-11, IFRS 17 & Directive 2009/138/CE, règlement délégué (UE) 2015/35 \\
\bottomrule
\end{tabular}
\end{table}
\begin{center}
\small Du \textbf{conservatisme comptable} au \textbf{réalisme économique}
\end{center}
\end{frame}

\begin{frame}{Un exemple : mêmes actifs, deux bilans}
\begin{columns}[T]
\begin{column}{0.52\textwidth}
\begin{table}
\small
\begin{tabular}{lrrr}
\toprule
 & \textbf{Comptable} & \textbf{Prudentiel} & \textbf{Écart} \\
\midrule
Actifs financiers & 1\,000 & 950 & $-50$ \\
Provisions techniques & 800 & 850 & $+50$ \\
Fonds propres & 200 & 100 & $-100$ \\
\bottomrule
\end{tabular}
\end{table}
\end{column}
\begin{column}{0.44\textwidth}
\begin{block}{Lecture}
\begin{itemize}
  \item Les actifs perdent leurs moins-values latentes
  \item Les engagements sont réévalués
  \item Les fonds propres sont \textbf{divisés par deux}
\end{itemize}
\end{block}
\end{column}
\end{columns}
\vspace{0.3cm}
\begin{alertblock}{Avec un SCR de 90}
Ratio de couverture $= 100 / 90 = 111\,\%$ : l'entreprise est couverte, avec peu de marge.
Les fonds propres prudentiels sont souvent plus faibles que les fonds propres comptables, mais
pas toujours : les plus-values latentes jouent dans l'autre sens.
\end{alertblock}
\end{frame}

\begin{frame}{Les acteurs du contrôle}
\begin{columns}[T]
\begin{column}{0.48\textwidth}
\begin{block}{EIOPA (en français AEAPP), niveau européen}
\small
\begin{itemize}
  \item Créée par le règlement (UE) n° 1094/2010, à Francfort
  \item Rédige orientations et projets de normes techniques
  \item Veille à la convergence des pratiques de contrôle
  \item Organise les tests de résistance européens
  \item Publie chaque mois la courbe des taux sans risque et l'ajustement symétrique
\end{itemize}
\end{block}
\end{column}
\begin{column}{0.48\textwidth}
\begin{block}{ACPR, niveau national}
\small
\begin{itemize}
  \item Autorité administrative indépendante adossée à la Banque de France
  \item Délivre les agréments
  \item Contrôle sur pièces (états remis via le portail e-Surfi) et sur place
  \item Prend des mesures de police administrative, sanctionne
  \item Prépare et conduit la résolution
\end{itemize}
\end{block}
\end{column}
\end{columns}
\end{frame}

\begin{frame}{Pourquoi contrôler les assureurs ?}
\begin{enumerate}
  \item \textbf{Protéger les assurés et les bénéficiaires} : c'est l'objectif principal du contrôle
        (directive 2009/138/CE, art.~27)
  \item \textbf{Préserver la stabilité financière} : les assureurs sont de grands investisseurs,
        leurs difficultés se propagent
  \item \textbf{Maintenir la confiance} : un engagement d'assurance vaut ce que vaut la promesse
        de l'assureur
\end{enumerate}
\vspace{0.4cm}
\begin{alertblock}{À retenir}
La comptabilité prudentielle est un outil de \textbf{contrôle de la solvabilité}, pas une mesure
de la performance.
\end{alertblock}
\end{frame}

% =====================================================================
\section{De Solvabilité I à Solvabilité II}
% =====================================================================

\begin{frame}{Solvabilité I : des marges forfaitaires}
\begin{columns}[T]
\begin{column}{0.48\textwidth}
\begin{block}{Le dispositif}
\begin{itemize}
  \item Premières directives : non-vie (1973), vie (1979), réformées en 2002
  \item Marge de solvabilité \textbf{forfaitaire} :
  \begin{itemize}
    \item vie : 4\,\% des provisions mathématiques $+$ 0,3\,\% des capitaux sous risque
    \item non-vie : environ 16 à 18\,\% des primes, ou 23 à 26\,\% des sinistres
  \end{itemize}
  \item Provisions et actifs pris dans les comptes sociaux
\end{itemize}
\end{block}
\end{column}
\begin{column}{0.48\textwidth}
\begin{alertblock}{Les limites}
\begin{itemize}
  \item Aucun lien entre exigence et \textbf{risques réellement pris} : deux assureurs aux
        placements opposés ont la même exigence
  \item Risques de marché, de crédit et opérationnels ignorés
  \item Règles nationales hétérogènes
\end{itemize}
\end{alertblock}
\end{column}
\end{columns}
\end{frame}

\begin{frame}{Solvabilité II : une exigence fondée sur les risques}
\begin{itemize}
  \item \textbf{Directive 2009/138/CE}, applicable depuis le \textbf{1\textsuperscript{er} janvier 2016}
  \item Un \textbf{bilan économique}, et un capital requis qui dépend du profil de risque de chaque entreprise
  \item La même logique que \textbf{Bâle II} pour les banques : capital fondé sur les risques, contrôle, discipline de marché
  \item Le principe de \textbf{proportionnalité} : exigences adaptées à la nature, à l'ampleur et à la complexité des risques
\end{itemize}
\vspace{0.3cm}
\begin{block}{Révision en cours}
La \textbf{directive (UE) 2025/2}, à appliquer à partir de 2027, révise notamment la marge de
risque, les mesures pour les engagements de long terme, la proportionnalité et l'information
publique. Les séances signalent ce qui change.
\end{block}
\end{frame}

\begin{frame}{Les trois piliers}
\begin{table}
\centering
\renewcommand{\arraystretch}{1.5}
\begin{tabular}{p{2.3cm}p{5.2cm}p{5.2cm}}
\toprule
\textbf{Pilier} & \textbf{Contenu} & \textbf{Ce qu'on y trouve} \\
\midrule
\textbf{1. Quantitatif} & Valorisation, fonds propres, capital requis & Bilan prudentiel, SCR, MCR \\
\textbf{2. Qualitatif} & Gouvernance, gestion des risques & Fonctions clés, ORSA \\
\textbf{3. Information} & Rapports au superviseur et au public & États QRT, RSR, SFCR \\
\bottomrule
\end{tabular}
\end{table}
\vspace{0.2cm}
\begin{center}
\small Cette séance couvre surtout le pilier 1 ; les piliers 2 et 3 sont présentés ici, l'ORSA
sera traité en séance 8.
\end{center}
\end{frame}

\begin{frame}{Pilier 2 : gouvernance et ORSA}
\begin{columns}[T]
\begin{column}{0.48\textwidth}
\begin{block}{Quatre fonctions clés}
\begin{itemize}
  \item Gestion des risques (art.~44)
  \item Vérification de la conformité (art.~46)
  \item Audit interne (art.~47)
  \item Fonction actuarielle (art.~48)
\end{itemize}
\end{block}
\end{column}
\begin{column}{0.48\textwidth}
\begin{block}{ORSA (art.~45)}
Évaluation interne des risques et de la solvabilité :
\begin{itemize}
  \item besoin global de solvabilité propre à l'entreprise
  \item respect permanent des exigences de capital
  \item écart entre son profil de risque et les hypothèses du SCR
\end{itemize}
Au moins une fois par an, et après tout changement important du profil de risque.
\end{block}
\end{column}
\end{columns}
\end{frame}

\begin{frame}{L'architecture des textes}
\begin{center}
\begin{tikzpicture}[
  niveau/.style={draw=bleuactu, fill=bleuactu!8, rounded corners, minimum width=10.5cm,
                 minimum height=0.85cm, align=center, font=\small},
  fleche/.style={-{Latex[length=2mm]}, thick, bleuactu}]
\node[niveau] (n1) {\textbf{Niveau 1} : directive 2009/138/CE --- les principes, transposés en droit national};
\node[niveau, below=0.35cm of n1] (n2) {\textbf{Niveau 2} : règlement délégué (UE) 2015/35 et normes techniques --- les règles de calcul, d'application directe};
\node[niveau, below=0.35cm of n2] (n3) {\textbf{Niveau 3} : orientations de l'EIOPA --- « se conformer ou s'expliquer »};
\node[niveau, below=0.35cm of n3, draw=vertactu, fill=vertactu!8] (fr) {\textbf{En France} : Code des assurances (transposition) ; instructions et notices de l'ACPR};
\draw[fleche] (n1) -- (n2);
\draw[fleche] (n2) -- (n3);
\draw[fleche] (n3) -- (fr);
\end{tikzpicture}
\end{center}
\small Les liens vers les textes sont sur le site du cours, page « Textes de référence ».
\end{frame}

% =====================================================================
\section{Le bilan prudentiel}
% =====================================================================

\begin{frame}{Le principe : la valeur d'échange}
\begin{block}{Directive 2009/138/CE, article 75}
\small
« Les actifs sont valorisés au montant pour lequel ils pourraient être échangés dans le cadre d'une
transaction conclue, dans des conditions de concurrence normales, entre des parties informées et
consentantes ; les passifs sont valorisés au montant pour lequel ils pourraient être transférés ou
réglés dans le cadre d'une transaction conclue, dans des conditions de concurrence normales, entre
des parties informées et consentantes. »
\end{block}
\begin{itemize}
  \item Une seule base pour tous les assureurs européens : comparabilité
  \item Le bilan ne suit ni les comptes sociaux ni les IFRS, même s'il s'en sert
  \item Les passifs ne tiennent \textbf{pas} compte de la qualité de crédit propre de l'entreprise
\end{itemize}
\end{frame}

\begin{frame}{Du bilan comptable au bilan prudentiel}
\begin{center}
\begin{tikzpicture}[
  bloc/.style={draw=white, very thick, minimum width=2.6cm, align=center, font=\scriptsize\bfseries, text=white},
  titre/.style={font=\small\bfseries, text=bleuactu}]
% Bilan comptable
\node[titre] at (2.6,4.6) {Bilan comptable};
\node[bloc, fill=bleuactu!70, minimum height=4.2cm] at (1.3,2.1) {Placements\\au coût\\historique};
\node[bloc, fill=rougeactu!70, minimum height=3.2cm] at (3.9,2.6) {Provisions\\techniques\\(prudentes)};
\node[bloc, fill=gray!70, minimum height=0.5cm] at (3.9,0.75) {Dettes};
\node[bloc, fill=vertactu!80, minimum height=0.5cm] at (3.9,0.25) {Capitaux propres};
% Flèche
\draw[-{Latex[length=3mm]}, ultra thick, orangeactu] (5.6,2.1) -- (7.0,2.1)
  node[midway, above, font=\scriptsize, text=orangeactu, align=center] {art.~75};
% Bilan prudentiel
\node[titre] at (9.9,4.6) {Bilan prudentiel};
\node[bloc, fill=bleuactu!70, minimum height=3.8cm] at (8.6,2.3) {Actifs en\\valeur de\\marché};
\node[bloc, fill=bleuactu!40, minimum height=0.4cm] at (8.6,0.2) {Impôts différés};
\node[bloc, fill=rougeactu!70, minimum height=2.6cm] at (11.2,2.9) {Meilleure\\estimation};
\node[bloc, fill=rougeactu!45, minimum height=0.5cm] at (11.2,1.35) {Marge de risque};
\node[bloc, fill=gray!70, minimum height=0.5cm] at (11.2,0.85) {Dettes};
\node[bloc, fill=vertactu!80, minimum height=0.6cm] at (11.2,0.3) {Excédent};
\end{tikzpicture}
\end{center}
\small
\textbf{Ce qui change} : plus- et moins-values latentes, provisions remplacées par meilleure estimation
$+$ marge de risque, écarts d'acquisition et frais d'acquisition reportés à zéro, impôts différés sur
tous les écarts. Schéma sans échelle.
\end{frame}

\begin{frame}{Valoriser les actifs}
\begin{block}{La hiérarchie de valorisation (règlement délégué, art.~10)}
\begin{enumerate}
  \item \textbf{Prix coté sur un marché actif} pour le même actif (action du CAC 40, obligation d'État)
  \item Sinon, \textbf{prix coté pour un actif similaire}, ajusté des différences
  \item Sinon, \textbf{méthode de valorisation alternative} (modèle), avec le plus possible de données de marché
\end{enumerate}
\end{block}
\begin{itemize}
  \item Les normes IFRS servent de référence quand elles sont cohérentes avec l'art.~75 (art.~9)
  \item Écart d'acquisition : valeur \textbf{nulle} ; autres incorporels : nuls sauf s'ils sont
        séparables et échangés sur un marché actif (art.~12)
  \item Exemple : obligation achetée 1\,000\,\euro, cotée 940\,\euro{} : ajustement de $-60$\,\euro
\end{itemize}
\end{frame}

\begin{frame}{Les provisions techniques}
\begin{center}
\fbox{\parbox{0.6\textwidth}{\centering\large Provisions techniques $=$ $\BE$ $+$ $\MR$}}
\end{center}
\vspace{0.2cm}
\begin{block}{Une valeur de transfert (art.~76)}
Le montant qu'il faudrait payer aujourd'hui à une autre entreprise pour qu'elle reprenne
\emph{sur-le-champ} les engagements d'assurance.
\end{block}
\begin{columns}[T]
\begin{column}{0.48\textwidth}
\begin{itemize}
  \item \textbf{Meilleure estimation} ($\BE$) : ce que coûteront en moyenne les engagements
\end{itemize}
\end{column}
\begin{column}{0.48\textwidth}
\begin{itemize}
  \item \textbf{Marge de risque} ($\MR$) : ce que demanderait le repreneur pour immobiliser le capital nécessaire
\end{itemize}
\end{column}
\end{columns}
\end{frame}

\begin{frame}{La meilleure estimation}
\begin{block}{Définition (art.~77)}
Moyenne pondérée par leur probabilité des flux de trésorerie futurs, actualisés avec la courbe
des taux sans risque pertinente :
\[ \BE = \sum_{t=1}^{n} \frac{\mathbb{E}[F_t]}{(1+r_t)^t} \]
\end{block}
\begin{columns}[T]
\begin{column}{0.48\textwidth}
\small
\textbf{Flux} : prestations, frais, primes futures dans les limites des contrats, participations
aux bénéfices futures. Calculée \textbf{brute} de réassurance.
\end{column}
\begin{column}{0.48\textwidth}
\small
\textbf{Hypothèses} réalistes : mortalité, rachats, frais, inflation, comportement des assurés.
\textbf{Taux} : courbe de l'EIOPA, mensuelle, extrapolée par Smith-Wilson.
\end{column}
\end{columns}
\vspace{0.2cm}
\begin{center}\small Calcul détaillé : séance 3 (non-vie) et séance 4 (vie).\end{center}
\end{frame}

\begin{frame}{La marge de risque : un coût du capital}
\begin{block}{Règlement délégué, art.~37 et 39}
\[ \MR = \mathrm{CoC} \times \sum_{t \geq 0} \frac{\SCR(t)}{(1+r_{t+1})^{t+1}}, \qquad \mathrm{CoC} = 6\,\% \]
$\SCR(t)$ : capital requis projeté pour l'entreprise de référence qui reprendrait les engagements,
jusqu'à leur extinction.
\end{block}
\begin{itemize}
  \item Logique : le repreneur doit immobiliser ce capital ; il demande une rémunération de 6\,\%
        par an au-delà du taux sans risque
  \item En pratique, le SCR futur est souvent approché par des simplifications (proportionnalité)
  \item La révision de Solvabilité II \textbf{allège} cette marge à partir de 2027
\end{itemize}
\end{frame}

\begin{frame}{Les impôts différés}
\begin{block}{Règlement délégué, art.~15}
Les impôts différés naissent des écarts entre la valeur prudentielle de chaque actif ou passif et
sa valeur fiscale.
\end{block}
\begin{columns}[T]
\begin{column}{0.48\textwidth}
\begin{itemize}
  \item Actifs réévalués à la hausse, passifs à la baisse : \textbf{impôt différé passif}
  \item Actifs réévalués à la baisse, passifs à la hausse : \textbf{impôt différé actif}
\end{itemize}
\end{column}
\begin{column}{0.48\textwidth}
\begin{alertblock}{Attention}
Un impôt différé actif n'est retenu que s'il est \textbf{recouvrable} : il faut des bénéfices
imposables futurs pour l'utiliser.
\end{alertblock}
\end{column}
\end{columns}
\vspace{0.3cm}
\small Exemple : 36\,000\keur{} de moins-values et de surcroît de provisions, taux d'impôt de 25\,\% :
impôt différé actif de 9\,000\keur.
\end{frame}

\begin{frame}{Des capitaux propres aux fonds propres prudentiels}
\begin{center}
\begin{tikzpicture}[
  etape/.style={draw=bleuactu, fill=bleuactu!8, rounded corners, minimum width=3.2cm,
                minimum height=1.2cm, align=center, font=\small},
  fleche/.style={-{Latex[length=2mm]}, thick, bleuactu}]
\node[etape] (a) {Excédent\\d'actif sur passif};
\node[etape, right=1.2cm of a] (b) {Fonds propres\\de base};
\node[etape, right=1.2cm of b] (c) {Fonds propres\\éligibles};
\draw[fleche] (a) -- node[above, font=\scriptsize]{$+$ dettes} node[below, font=\scriptsize]{subordonnées} (b);
\draw[fleche] (b) -- node[above, font=\scriptsize]{niveaux et} node[below, font=\scriptsize]{limites} (c);
\end{tikzpicture}
\end{center}
\begin{itemize}
  \item \textbf{Fonds propres de base} (art.~88) : excédent d'actif sur passif, plus les passifs
        subordonnés, moins les actions propres
  \item \textbf{Fonds propres auxiliaires} (art.~89--90) : éléments non versés mais appelables
        (capital non appelé, lettres de crédit), sur autorisation du superviseur
  \item \textbf{Réserve de réconciliation} (règlement délégué, art.~70) : l'excédent, diminué des
        autres éléments de fonds propres et des dividendes prévisibles. Elle porte tous les écarts de valorisation
\end{itemize}
\end{frame}

\begin{frame}{Trois niveaux de qualité}
\begin{table}
\centering
\small
\renewcommand{\arraystretch}{1.4}
\begin{tabular}{p{1.6cm}p{6.2cm}p{5.2cm}}
\toprule
\textbf{Niveau} & \textbf{Exemples} & \textbf{Qualité} \\
\midrule
\textbf{1} & Capital social versé, primes d'émission, réserve de réconciliation ; « restreint » : certains instruments subordonnés perpétuels & Disponible en permanence, totalement subordonné \\
\textbf{2} & Dettes subordonnées d'échéance initiale d'au moins 10 ans ; la plupart des fonds propres auxiliaires & Subordonné, durée suffisante \\
\textbf{3} & Impôts différés actifs nets ; dettes subordonnées d'au moins 5 ans & Qualité la plus faible \\
\bottomrule
\end{tabular}
\end{table}
\small Critères : directive, art.~93--94 ; listes : règlement délégué, art.~69 à 78. Les fonds propres
auxiliaires ne sont jamais de niveau 1.
\end{frame}

\begin{frame}{Les limites d'éligibilité}
\begin{columns}[T]
\begin{column}{0.48\textwidth}
\begin{block}{Pour couvrir le SCR}
\begin{itemize}
  \item Niveau 1 : au moins 50\,\% du SCR
  \item Niveau 3 : moins de 15\,\% du SCR
  \item Niveaux 2 $+$ 3 : au plus 50\,\% du SCR
  \item Niveau 1 restreint : au plus 20\,\% du niveau 1
\end{itemize}
\end{block}
\end{column}
\begin{column}{0.48\textwidth}
\begin{block}{Pour couvrir le MCR}
\begin{itemize}
  \item Niveau 1 : au moins 80\,\% du MCR
  \item Niveau 2 : au plus 20\,\% du MCR
  \item Niveau 3 : exclu
\end{itemize}
\end{block}
\end{column}
\end{columns}
\vspace{0.3cm}
\small Règlement délégué, art.~82. Les fonds propres au-delà des limites existent, mais ne
comptent pas dans la couverture.
\end{frame}

% =====================================================================
\section{Capital requis et ratio de couverture}
% =====================================================================

\begin{frame}{SCR et MCR}
\begin{columns}[T]
\begin{column}{0.48\textwidth}
\begin{block}{SCR (art.~100--101)}
\begin{itemize}
  \item Valeur en risque des fonds propres de base à \textbf{99,5\,\%}, horizon \textbf{un an} :
        résister à un choc « une fois en 200 ans »
  \item Formule standard (modules de risque agrégés) ou modèle interne approuvé
  \item Calculé au moins une fois par an
\end{itemize}
\end{block}
\end{column}
\begin{column}{0.48\textwidth}
\begin{block}{MCR (art.~128--129)}
\begin{itemize}
  \item Calibré sur une valeur en risque à 85\,\%, un an
  \item Formule linéaire, encadrée entre \textbf{25\,\%} et \textbf{45\,\%} du SCR, avec un plancher absolu
  \item Le seuil en deçà duquel l'agrément est en jeu
\end{itemize}
\end{block}
\end{column}
\end{columns}
\vspace{0.2cm}
\begin{center}\small Calcul de la formule standard : séance 7.\end{center}
\end{frame}

\begin{frame}{Le ratio de couverture et l'échelle d'intervention}
\begin{center}
\fbox{\parbox{0.6\textwidth}{\centering Ratio de couverture $=$ fonds propres éligibles $/$ SCR}}
\end{center}
\begin{table}
\centering
\small
\renewcommand{\arraystretch}{1.4}
\begin{tabular}{p{4.2cm}p{8.3cm}}
\toprule
\textbf{Situation} & \textbf{Conséquence} \\
\midrule
Ratio $\geq 100\,\%$ & Situation conforme ; l'ORSA dit quelle marge viser au-delà \\
SCR non couvert & Information immédiate de l'ACPR ; \textbf{plan de rétablissement} sous deux mois ;
retour à la couverture en six mois, prolongeable (art.~138) \\
MCR non couvert & \textbf{Plan de financement à court terme} sous un mois ; retour à la couverture en
trois mois (art.~139) ; à défaut, \textbf{retrait d'agrément} (art.~144) \\
\bottomrule
\end{tabular}
\end{table}
\end{frame}

\begin{frame}{Exemple : l'Assureur Alpha}
\begin{columns}[T]
\begin{column}{0.44\textwidth}
\begin{block}{Situation initiale}
\begin{itemize}
  \item Fonds propres éligibles : 120\,M\euro
  \item SCR : 100\,M\euro{} ; MCR : 45\,M\euro
  \item Ratio : $120/100 = 120\,\%$
\end{itemize}
\end{block}
\end{column}
\begin{column}{0.52\textwidth}
\begin{alertblock}{Après un choc de $-20$\,M\euro{} sur les actifs}
\begin{itemize}
  \item Fonds propres : 100\,M\euro{}, ratio : $100\,\%$
  \item L'entreprise est exactement au seuil : le moindre choc supplémentaire déclenche
        l'intervention du superviseur
\end{itemize}
\end{alertblock}
\end{column}
\end{columns}
\vspace{0.3cm}
\small \textbf{Simplifications} : le SCR ne change pas après le choc, et la perte n'est pas allégée
par les impôts différés. Deux effets que l'on mesurera dans les séances suivantes.
\end{frame}

% =====================================================================
\section{Rendre compte : le pilier 3}
% =====================================================================

\begin{frame}{Les états quantitatifs (QRT)}
\begin{table}
\centering
\small
\renewcommand{\arraystretch}{1.3}
\begin{tabular}{lll}
\toprule
\textbf{État} & \textbf{Contenu} & \textbf{Séance} \\
\midrule
S.02.01 & Bilan prudentiel & 1 \\
S.06.02 & Liste des actifs, ligne à ligne & 2 \\
S.12.01 et S.17.01 & Provisions techniques vie et non-vie & 3, 4 \\
S.23.01 & Fonds propres & 1, 7 \\
S.25.01 & SCR en formule standard & 7 \\
S.28.01 & MCR & 7 \\
\bottomrule
\end{tabular}
\end{table}
\begin{itemize}
  \item Modèles fixés par le règlement d'exécution (UE) 2023/894 ; remise annuelle et trimestrielle
  \item Format \textbf{XBRL}, taxonomie de l'EIOPA ; remis à l'ACPR, qui les transmet à l'EIOPA
\end{itemize}
\end{frame}

\begin{frame}{Les rapports narratifs : SFCR et RSR}
\begin{columns}[T]
\begin{column}{0.48\textwidth}
\begin{block}{SFCR : rapport sur la solvabilité et la situation financière}
\small
\textbf{Public}, annuel, sur le site de l'entreprise (directive, art.~51). Cinq parties :
\begin{enumerate}[A.]
  \item Activité et résultats
  \item Système de gouvernance
  \item Profil de risque
  \item Valorisation à des fins de solvabilité
  \item Gestion du capital
\end{enumerate}
avec une sélection d'états en annexe.
\end{block}
\end{column}
\begin{column}{0.48\textwidth}
\begin{block}{RSR : rapport régulier au contrôleur}
\small
\textbf{Non public}, destiné à l'ACPR (directive, art.~35). Même plan, plus détaillé :
informations sensibles, stratégie, réassurance, limites de risque.
\end{block}
\begin{alertblock}{Qui produit quoi ?}
\small
Calculs : actuariat et direction financière ; validation des provisions : fonction actuarielle ;
signature : dirigeants effectifs ; approbation : conseil d'administration.
\end{alertblock}
\end{column}
\end{columns}
\end{frame}

% =====================================================================
\section{Étude de cas : le mini-bilan prudentiel}
% =====================================================================

\begin{frame}{L'Assureur Bêta : le bilan des comptes sociaux}
\begin{table}
\centering
\small
\begin{tabular}{lrr@{\qquad}lr}
\toprule
\textbf{Actif} & \textbf{Comptable} & \textbf{Marché} & \textbf{Passif} & \textbf{Comptable} \\
\midrule
Obligations d'État        & 500\,000 & 480\,000 & Provisions techniques       & 950\,000 \\
Obligations d'entreprises & 200\,000 & 195\,000 & Dettes subordonnées (12 ans) & 50\,000 \\
Actions cotées            & 100\,000 &  90\,000 & Autres dettes               &  20\,000 \\
Immobilier                & 200\,000 & 210\,000 & Capital social              &  30\,000 \\
Trésorerie                &  50\,000 &  50\,000 & Réserves et résultat        &  20\,000 \\
Créances                  &  20\,000 &  19\,000 &                             &          \\
\midrule
\textbf{Total}            & \textbf{1\,070\,000} & & \textbf{Total} & \textbf{1\,070\,000} \\
\bottomrule
\end{tabular}
\end{table}
\small Assureur vie fictif, montants en k\euro. Meilleure estimation : 920\,000\keur ; taux d'impôt
25\,\% ; SCR 70\,000\keur, MCR 17\,500\keur. Données : \texttt{assureur-beta.csv} sur le site du cours.
\end{frame}

\begin{frame}{La démarche du TD, en sept étapes}
\begin{enumerate}
  \item \textbf{Revaloriser les actifs} en valeur de marché
  \item \textbf{Recalculer les provisions techniques} : meilleure estimation $+$ marge de risque
  \item \textbf{Calculer les impôts différés} sur les écarts de valorisation
  \item \textbf{Construire le bilan prudentiel} et l'excédent d'actif sur passif ; le réconcilier avec les capitaux propres
  \item \textbf{Classer les fonds propres} en niveaux ; calculer la réserve de réconciliation
  \item \textbf{Vérifier la couverture} du SCR et du MCR après limites d'éligibilité
  \item \textbf{Mesurer la sensibilité} à trois chocs : actions, taux, écarts de crédit
\end{enumerate}
\vspace{0.2cm}
\begin{alertblock}{Avant de calculer}
Pariez : les fonds propres prudentiels de Bêta seront-ils supérieurs ou inférieurs à ses
50\,000\keur{} de capitaux propres comptables ? Son SCR sera-t-il couvert ?
\end{alertblock}
\end{frame}

% =====================================================================
\section{Synthèse}
% =====================================================================

\begin{frame}{À retenir}
\begin{enumerate}
  \item La comptabilité prudentielle répond à une question : l'assureur pourra-t-il payer, même
        après un choc ? Elle ne mesure pas sa performance.
  \item Le bilan prudentiel valorise tout à la \textbf{valeur d'échange} (art.~75) : actifs en valeur
        de marché, provisions $=$ meilleure estimation $+$ marge de risque, impôts différés sur les écarts.
  \item Les fonds propres se \textbf{classent} en trois niveaux et ne comptent que dans des
        \textbf{limites} (art.~82).
  \item Le \textbf{ratio de couverture} rapporte les fonds propres éligibles au SCR (99,5\,\%, un an) ;
        sous le SCR puis sous le MCR, le superviseur intervient de plus en plus fort.
  \item Trois piliers : le chiffre (1), la gouvernance (2), la transparence (3).
\end{enumerate}
\end{frame}

\begin{frame}{Pour la prochaine séance}
\begin{block}{À faire}
\begin{itemize}
  \item Terminer le TD 1 ; le corrigé est publié sur le site après la séance
  \item Revoir les unités 1 à 3 du \textbf{polycopié de prérequis} : valeur actuelle, échéancier,
        obligations et courbe des taux
\end{itemize}
\end{block}
\begin{block}{Questions de rappel, en ouverture de la séance 2}
\begin{itemize}
  \item Pourquoi un écart d'acquisition vaut-il zéro au bilan prudentiel ?
  \item D'où vient un impôt différé actif, et à quelle condition est-il retenu ?
  \item Pourquoi un SCR non couvert n'entraîne-t-il pas immédiatement un retrait d'agrément ?
\end{itemize}
\end{block}
\end{frame}

\begin{frame}[c]
\centering
\Huge Questions ?

\vspace{1cm}
\Large
\textbf{Séance 2} : valorisation des produits financiers
\end{frame}

\end{document}
